\documentclass[11pt]{article}
\usepackage[margin=1in]{geometry}
\usepackage[T1]{fontenc}
\usepackage{lmodern}
\usepackage{microtype}
\usepackage{booktabs}
\usepackage{tabularx}
\usepackage{amsmath}
\usepackage{xcolor}
\usepackage[hidelinks]{hyperref}
\usepackage{enumitem}
\setlist{nosep}

% Draft markers: every \todo must be gone before arXiv submission.
\newcommand{\todo}[1]{\textcolor{red}{[TODO: #1]}}

\title{Recovery That Re-Asks the Model:\\
Duplicate and Lost Side Effects in Crash-Recovered LLM Agents}
\author{Param Madan\\Northeastern University\\\texttt{madan.pa@northeastern.edu}}
\date{\todo{date}}

\begin{document}
\maketitle

\begin{abstract}
Durable-execution systems make crash recovery correct by replaying
deterministic code against recorded results. LLM agents break this premise:
the code that decides which side effects to perform is a model whose answer
is not reproducible across requests. We show that recovery schemes that
\emph{re-ask the model} for a decision the crashed process had already
received can execute side effects twice or, depending on the
idempotency-key scheme, silently drop them. We separate two design axes that
are usually conflated, \emph{recording granularity} (is the model's decision
durable before it is acted on?) and \emph{key derivation} (provider tool-call
id, argument hash, or position), and measure their interaction under a
physical after-effect crash in two pre-registered rounds
\todo{N runs, task shapes, results sentence}. A survey of the official
documentation of 15 agent-framework configurations finds
\todo{survey sentence}. We argue that recording the model's decision before
acting is a correctness requirement, not an optimization, and outline a
correctness condition for recovered agent runs.
\end{abstract}

\section{Introduction}

A customer-service agent looks up a case, decides to issue a \$25 service
credit, and calls the ledger. The ledger writes the credit. Before the
response reaches the agent, its process is killed. A recovery process starts
the turn over: it sends the same conversation to the model, the model decides
again to issue the credit, and the provider assigns the new tool call a new
identifier. If that identifier is the idempotency key, the ledger now holds
two credits.

Idempotency keys are the standard defense against retried side effects, but
a key identifies a \emph{decision}. When recovery causes the decision to be
made again, the key is made again too, and whether it matches depends on
whether the new decision is identical to the old one \emph{in whatever the key
scheme looks at}. Provider call ids are fresh on every request; argument
hashes change whenever the model writes any argument in its own words;
position-based keys break when the model changes the number or order of its
effects, and in that case they fail in the other direction, by suppressing a
legitimate effect.

\paragraph{Contributions.}
\begin{enumerate}
  \item A taxonomy of agent recovery along two axes, recording granularity and
        key derivation, with the condition under which each pair is
        exactly-once (\S\ref{sec:taxonomy}).
  \item Two pre-registered measurement rounds under a physical after-effect
        crash, with predictions committed before data (\S\ref{sec:method},
        \S\ref{sec:results}).
  \item A survey of what agent frameworks' official documentation prescribes
        (\S\ref{sec:survey}).
  \item An open harness, raw per-run data, and both pre-registrations.
        \todo{repository URL once public}
\end{enumerate}

\section{Background}
\todo{Durable execution and deterministic replay; sagas and compensation;
idempotency keys and exactly-once processing; tool-calling agents and
provider-assigned call ids. Citations pending LibKey resolution.}

\section{A taxonomy of agent recovery}
\label{sec:taxonomy}

\paragraph{Recording granularity.} Under \emph{step} recording, every model
response, including the tool calls it chose, is durable before any of those
calls execute, and every tool result is durable when it returns. Recovery
reuses recorded responses and re-executes only calls whose results are
missing. Under \emph{turn} recording, state is persisted when a turn or run
completes; a crash inside the turn loses the model's decision, and recovery
asks the model again.

\paragraph{Key derivation.} \texttt{tool\_call\_id} uses the identifier the
model provider assigns to the call. \texttt{args\_hash} hashes the task id,
tool name, and canonicalized arguments. \texttt{position} combines the task
id, tool name, and the ordinal of that tool's call within the run.

\begin{table}[t]
\centering
\begin{tabularx}{\linewidth}{lXXX}
\toprule
 & \texttt{tool\_call\_id} & \texttt{args\_hash} & \texttt{position} \\
\midrule
step & safe & safe & safe \\
turn & unsafe: a new id per request & unsafe when the model authors any argument
     & unsafe when the number or order of effects can change; drops effects \\
\bottomrule
\end{tabularx}
\caption{Predicted safety of each granularity/key pair under an after-effect
crash. \S\ref{sec:results} tests every cell.}
\label{tab:taxonomy}
\end{table}

\paragraph{Why step recording makes the key irrelevant.} Under step recording
the call being retried is the \emph{same recorded call}, so any deterministic
function of it (its provider id, its arguments, or its position) yields the
same key on retry. Under turn recording the retried call is a \emph{new}
decision, and exactly-once holds only when the new decision agrees with the
old one on every input to the key function.

\section{Method}
\label{sec:method}

\paragraph{Harness.} The tool server owns the ledger (ground truth), which
deduplicates by idempotency key. The after-effect fault: the server writes
the credit, then holds the response; the harness kills the agent process
with \texttt{SIGKILL}, releases the hold, and runs the agent's recovery (the
same command) once. Scoring reads only the ledger and an append-only
observation trace the agent never reads.

\paragraph{Study agent.} One neutral, roughly 200-line tool-calling agent
implements both granularities and all three key schemes, so the comparison is
between recording disciplines rather than products. The author also builds a
durable runtime (Aurelia) that uses step recording; that runtime is
deliberately \emph{not} a contestant.

\paragraph{Tasks.} \emph{credit}: one case, every tool argument copied from a
tool result. \emph{memo}: as credit, plus a required customer-facing reason
the model writes. \emph{two}: two cases in one request, order chosen by the
model.

\paragraph{Models and sampling.} \texttt{qwen2.5:3b} and
\texttt{llama3.2:3b} served by Ollama~0.32.5 at temperatures 0, 0.7, and 1.0.
The per-request seed differs between the crashed and the recovering
invocation, as in production, and is reproducible from the run seed.

\paragraph{Pre-registration.} Round~1 (credit, Apple~M1, serial; 1,080 runs)
and Round~2 (all three tasks, one NVIDIA A100, 16 concurrent runs; 3,240
runs, including a hardware replication of Round~1) were each pre-registered
with predictions, outcomes, exclusions, and analysis before any measured run.
The Round~1 predictions were drafted with an AI coding assistant and approved
unedited; we report every prediction, including falsified ones.

\paragraph{Outcomes.} \emph{Duplicate}: more than one ledger credit for a
case. \emph{Lost effect}: the run reports success but the ledger lacks a
credit for a named case. Rates are over runs in which the kill fired; runs in
which it did not are reported per cell. Intervals are Wilson 95\%.

\section{Results}
\label{sec:results}
\todo{Round 1 table and hypothesis verdicts from results-r1.md.}

\todo{Round 2 tables and verdicts from results-r2.md, falsified predictions
first.}

\section{What frameworks document}
\label{sec:survey}
\todo{From rq4-survey.md: step for the durable-execution engines and their
integrations; turn for agent-SDK session stores and Microsoft's Durable Task
extension for Agent Framework (inferred from samples; re-verify wording);
quickstarts mostly none or turn; no framework names tool\_call\_id as a key.
Access date 2026-10-12.}

\section{Discussion}
\todo{Implications for framework defaults and docs; cost of step recording;
towards a correctness condition: a recovered run is correct iff its effect
history equals that of some crash-free run of the same agent.}

\section{Threats to validity}
Small open models; one domain; we implement the patterns frameworks
document rather than running every framework's code; Round~2 uses different
hardware and concurrent decoding (hence the replication task); the author
builds a step-recording runtime. Code, configuration, raw data, and both
pre-registrations are public.

\section{Related work}
\todo{Pending LibKey: durable-execution systems, workflow replay, sagas,
exactly-once stream processing, agent reliability and fault-injection
benchmarks.}

\end{document}
